\documentclass[../../../../main.tex]{subfiles}
\begin{document}

\begin{flushright}
\footnotesize\textit{【自動文字起こし・要確認】}
\end{flushright}

\begin{proof}[解]
$P(x)$ の $n$ 次以下の部分を $R_n(x)$, $Q_n(x) = (1+x)^k R_n(x)$ とおくと
「$Q_n(x)$ の $n$ 次以下の項の係数が $\mathbb{Z}$ に属する時, $R_n(x)$ の全ての項の係数も $\mathbb{Z}$ に属する」ことを示せば良い.
\[
R_n(x) = \sum_{i=0}^n a_i x^i
\]
\[
Q_n(x) = \sum_{i=0}^{n+k} b_i x^i \quad (b_i \text{ は } i=0,1,\dots,n+k \text{ の時整数}) \quad \cdots \star
\]
以下, $\star$ を $l$ についての帰納法で示す. $l=0$ の時,
$b_0 = a_0$ となり, $a_0 \in \mathbb{Z}$ だから成立する. 以下 $i \le l$ ($l=1,2,\dots,n-1$) での成立を仮定する.
\[
Q_n(x) = \left( x^k + k x^{k-1} + {}_k\mathrm{C}_2 x^{k-2} + \dots + kx + 1 \right) (a_n x^n + \dots + a_0)
\]
だから, $Q_n(x)$ の $l+1$ 次項は, $a_{l+1} + (\text{整数}) \quad (\because \text{仮定})$ の形で表せるから
\[
a_{l+1} = b_{l+1} - (\text{整数}) \in \mathbb{Z}
\]
となり, $i=l+1$ でも仮定は成立. 以上から示された. \quad \text{//}
\end{proof}

\end{document}
